Bagian ini membahas hasil percobaan algoritma pathfinding yang telah dilakukan pada lima test case, yaitu Test 8, Test 17, Test 20, Test 24, dan Test 25. 
Analisis dilakukan berdasarkan beberapa metrik utama, yaitu total cost solusi, waktu eksekusi, jumlah iterasi, stabilitas algoritma, dan pengaruh heuristik terhadap proses pencarian.

\begin{tcolorbox}[
	colback=blue!5,
	colframe=blue!60!black,
	title=\textbf{Tujuan Analisis Percobaan},
	arc=2mm,
	boxrule=0.8pt
]
Analisis ini bertujuan untuk membandingkan perilaku UCS, GBFS, A*, dan Weighted A* pada berbagai bentuk map. 
Setiap algoritma diuji untuk melihat kemampuan menemukan solusi, efisiensi jumlah iterasi, kecepatan eksekusi, serta kualitas cost solusi yang dihasilkan.
\end{tcolorbox}


\section{Deskripsi Umum Test Case}

Kelima test case yang digunakan memiliki karakteristik yang berbeda. 
Perbedaan tersebut meliputi ukuran papan, jumlah angka wajib, keberadaan lava, banyaknya rintangan, serta variasi cost tile. 
Perbedaan karakteristik ini penting karena dapat memengaruhi jumlah state yang dieksplorasi oleh setiap algoritma.

\begin{table}[H]
	\centering
	\caption{Deskripsi umum test case}
	\small
	\begin{tabular}{|p{2cm}|p{2cm}|p{2.2cm}|p{2cm}|p{4.8cm}|}
		\hline
		\rowcolor{blue!15}
		\textbf{Test Case} & \textbf{Ukuran} & \textbf{Angka Wajib} & \textbf{Lava} & \textbf{Karakteristik Map} \\
		\hline
		Test 8 & $6 \times 6$ & Tidak ada & Tidak ada & Map kecil dan sederhana untuk menguji mekanisme sliding dasar. \\
		\hline
		Test 17 & $10 \times 10$ & Ada, 0--7 & Tidak ada & Map besar dengan banyak angka dan variasi cost tile. \\
		\hline
		Test 20 & $12 \times 12$ & Tidak ada & Tidak ada & Map berbentuk maze dengan banyak rintangan. \\
		\hline
		Test 24 & $7 \times 12$ & Ada, 0--6 & Ada & Map dengan kombinasi angka berurutan, obstacle, dan lava. \\
		\hline
		Test 25 & $12 \times 12$ & Tidak ada & Ada & Map besar dengan banyak lava dan obstacle. \\
		\hline
	\end{tabular}
\end{table}

\begin{tcolorbox}[
	colback=gray!5,
	colframe=black!60,
	title=\textbf{Klasifikasi Tingkat Kesulitan Test Case},
	arc=2mm,
	boxrule=0.8pt
]
\begin{itemize}
	\item \textbf{Test 8} merupakan kasus sederhana karena ukuran kecil dan tidak memiliki angka maupun lava.
	\item \textbf{Test 17} menantang karena memiliki banyak angka yang harus dilewati secara berurutan.
	\item \textbf{Test 20} menantang dari sisi rintangan karena bentuknya menyerupai maze.
	\item \textbf{Test 24} menantang karena menggabungkan angka berurutan dan lava.
	\item \textbf{Test 25} menantang karena memiliki banyak lava dan obstacle, meskipun tidak memiliki angka wajib.
\end{itemize}
\end{tcolorbox}


\section{Metrik yang Dibandingkan}

Pada setiap percobaan, beberapa metrik digunakan untuk membandingkan performa algoritma. 
Metrik tersebut tidak hanya melihat apakah solusi ditemukan, tetapi juga memperhatikan efisiensi dan kualitas solusi.

\begin{table}[H]
	\centering
	\caption{Metrik pembanding hasil percobaan}
	\begin{tabular}{|p{4cm}|p{10cm}|}
		\hline
		\rowcolor{gray!20}
		\textbf{Metrik} & \textbf{Makna} \\
		\hline
		Total cost & Menunjukkan biaya total solusi yang ditemukan. Semakin kecil nilainya, semakin baik solusi dari sisi optimalitas cost. \\
		\hline
		Waktu eksekusi & Menunjukkan waktu yang dibutuhkan algoritma untuk menemukan solusi dalam milidetik. \\
		\hline
		Jumlah iterasi & Menunjukkan banyaknya state atau konfigurasi yang ditinjau selama pencarian. \\
		\hline
		Stabilitas & Menunjukkan kemampuan algoritma menghasilkan performa yang konsisten pada berbagai bentuk map. \\
		\hline
		Optimalitas & Menunjukkan apakah algoritma menghasilkan cost solusi minimum dibanding algoritma lain. \\
		\hline
		Efisiensi pencarian & Menunjukkan kemampuan algoritma menemukan solusi dengan jumlah iterasi dan waktu yang rendah. \\
		\hline
	\end{tabular}
\end{table}


\section{Rekapitulasi Hasil Percobaan}

Tabel berikut merangkum hasil terbaik dari setiap test case berdasarkan cost solusi, waktu eksekusi, dan jumlah iterasi.

\begin{longtable}{|p{2cm}|p{3.2cm}|p{3.2cm}|p{3.2cm}|}
	\caption{Rekapitulasi hasil terbaik per test case} \\
	\hline
	\rowcolor{green!15}
	\textbf{Test Case} & \textbf{Cost Terbaik} & \textbf{Waktu Tercepat} & \textbf{Iterasi Tersedikit} \\
	\hline
	\endfirsthead
	
	\hline
	\rowcolor{green!15}
	\textbf{Test Case} & \textbf{Cost Terbaik} & \textbf{Waktu Tercepat} & \textbf{Iterasi Tersedikit} \\
	\hline
	\endhead
	
	Test 8 & Semua algoritma, cost 37 & GBFS - Basic, 0.077 ms & GBFS, 3 iterasi \\
	\hline
	Test 17 & Semua algoritma, cost 1085 & GBFS - Custom, 0.439 ms & GBFS - Advanced, 10 iterasi \\
	\hline
	Test 20 & UCS, A*, WA*, cost 810 & GBFS - Basic, 0.065 ms & GBFS, 3 iterasi \\
	\hline
	Test 24 & Semua algoritma, cost 719 & GBFS - Custom, 0.192 ms & GBFS - Custom dan Advanced, 7 iterasi \\
	\hline
	Test 25 & Semua algoritma, cost 1637 & UCS, 0.159 ms & UCS dan GBFS, 11 iterasi \\
	\hline
\end{longtable}

\begin{tcolorbox}[
	colback=green!5,
	colframe=green!50!black,
	title=\textbf{Gambaran Umum Hasil},
	arc=2mm,
	boxrule=0.8pt
]
Sebagian besar test case menunjukkan bahwa semua algoritma dapat menemukan solusi valid. 
Perbedaan utama terlihat pada jumlah iterasi dan waktu eksekusi. 
Namun, pada Test 20, GBFS menghasilkan cost yang jauh lebih besar daripada UCS, A*, dan Weighted A*, sehingga terlihat bahwa GBFS tidak selalu aman untuk mencari solusi dengan cost minimum.
\end{tcolorbox}


\section{Analisis Per Test Case}

\subsection{Analisis Test Case 8}

Test Case 8 merupakan map kecil dan sederhana. 
Pada test case ini, semua algoritma menghasilkan cost solusi yang sama, yaitu 37. 
Hal ini menunjukkan bahwa pada map kecil tanpa angka dan lava, ruang pencarian relatif sederhana sehingga perbedaan strategi algoritma tidak terlalu memengaruhi kualitas cost solusi.

\begin{table}[H]
	\centering
	\caption{Analisis ringkas Test Case 8}
	\begin{tabular}{|p{4cm}|p{10cm}|}
		\hline
		\rowcolor{blue!15}
		\textbf{Aspek} & \textbf{Analisis} \\
		\hline
		Cost solusi & Semua algoritma menghasilkan cost 37. \\
		\hline
		Jumlah iterasi & GBFS membutuhkan 3 iterasi, lebih sedikit dibanding UCS yang membutuhkan 4 iterasi dan A*/WA* yang membutuhkan 5 iterasi. \\
		\hline
		Waktu eksekusi & Perbedaan waktu tidak terlalu signifikan karena ukuran map kecil. \\
		\hline
		Pengaruh heuristik & Tidak terlalu besar karena tidak ada angka dan struktur map sederhana. \\
		\hline
		Kesimpulan & Semua algoritma efektif pada map kecil, tetapi GBFS lebih efisien dari sisi iterasi. \\
		\hline
	\end{tabular}
\end{table}

\subsection{Analisis Test Case 17}

Test Case 17 memiliki banyak angka yang harus dilewati secara berurutan. 
Pada kasus ini, semua algoritma tetap menghasilkan cost solusi yang sama, yaitu 1085. 
Namun, jumlah iterasi berbeda cukup signifikan antara algoritma berbasis heuristik dan UCS/A*.

\begin{table}[H]
	\centering
	\caption{Analisis ringkas Test Case 17}
	\begin{tabular}{|p{4cm}|p{10cm}|}
		\hline
		\rowcolor{blue!15}
		\textbf{Aspek} & \textbf{Analisis} \\
		\hline
		Cost solusi & Semua algoritma menghasilkan cost 1085. \\
		\hline
		Constraint angka & Banyak angka menyebabkan state harus menyimpan progres angka melalui collected mask. \\
		\hline
		Jumlah iterasi & UCS dan A* berada di sekitar 71--72 iterasi, sedangkan GBFS dapat turun hingga 10--25 iterasi. \\
		\hline
		Pengaruh heuristik & Heuristik sangat membantu GBFS karena mampu mengarahkan pencarian menuju target angka berikutnya. \\
		\hline
		Kesimpulan & Pada map dengan banyak angka, heuristik yang sesuai dapat mengurangi jumlah iterasi secara signifikan. \\
		\hline
	\end{tabular}
\end{table}

\subsection{Analisis Test Case 20}

Test Case 20 merupakan map berbentuk maze dengan banyak rintangan. 
Pada test case ini, UCS, A*, dan Weighted A* menghasilkan cost 810, sedangkan GBFS menghasilkan cost 6174. 
Perbedaan ini menunjukkan risiko GBFS ketika heuristik terlalu berfokus pada kedekatan posisi tanpa mempertimbangkan cost aktual.

\begin{table}[H]
	\centering
	\caption{Analisis ringkas Test Case 20}
	\begin{tabular}{|p{4cm}|p{10cm}|}
		\hline
		\rowcolor{orange!20}
		\textbf{Aspek} & \textbf{Analisis} \\
		\hline
		Cost solusi & UCS, A*, dan WA* menghasilkan cost 810, sedangkan GBFS menghasilkan cost 6174. \\
		\hline
		Banyak rintangan & Rintangan membuat heuristik jarak menjadi kurang akurat karena jalur yang tampak dekat belum tentu murah. \\
		\hline
		Jumlah iterasi & GBFS hanya membutuhkan 3 iterasi, tetapi menghasilkan cost yang jauh lebih besar. \\
		\hline
		Optimalitas & UCS, A*, dan WA* lebih baik dari sisi cost pada test case ini. \\
		\hline
		Kesimpulan & GBFS cepat, tetapi pada map maze dapat memilih jalur yang tidak optimal dari sisi cost. \\
		\hline
	\end{tabular}
\end{table}

\subsection{Analisis Test Case 24}

Test Case 24 menggabungkan angka berurutan dan lava. 
Semua algoritma menghasilkan cost yang sama, yaitu 719. 
Meskipun terdapat lava, program tetap dapat menyaring jalur tidak valid melalui mekanisme \texttt{MovementEngine}.

\begin{table}[H]
	\centering
	\caption{Analisis ringkas Test Case 24}
	\begin{tabular}{|p{4cm}|p{10cm}|}
		\hline
		\rowcolor{blue!15}
		\textbf{Aspek} & \textbf{Analisis} \\
		\hline
		Cost solusi & Semua algoritma menghasilkan cost 719. \\
		\hline
		Lava & Gerakan yang melewati lava dianggap tidak valid sehingga tidak masuk ke state berikutnya. \\
		\hline
		Angka berurutan & Algoritma harus tetap mengikuti urutan angka sebelum menuju goal. \\
		\hline
		Jumlah iterasi & GBFS dengan Custom dan Advanced hanya membutuhkan 7 iterasi, lebih sedikit dibanding UCS, A*, dan WA*. \\
		\hline
		Kesimpulan & Heuristik efektif pada map ini karena mampu mengarahkan pencarian meskipun terdapat lava dan angka wajib. \\
		\hline
	\end{tabular}
\end{table}

\subsection{Analisis Test Case 25}

Test Case 25 memiliki banyak lava dan obstacle, tetapi tidak memiliki angka wajib. 
Semua algoritma menghasilkan cost yang sama, yaitu 1637. 
Hal ini menunjukkan bahwa meskipun map cukup kompleks, jalur solusi yang valid relatif terarah sehingga semua algoritma dapat menemukan solusi dengan cost yang sama.

\begin{table}[H]
	\centering
	\caption{Analisis ringkas Test Case 25}
	\begin{tabular}{|p{4cm}|p{10cm}|}
		\hline
		\rowcolor{blue!15}
		\textbf{Aspek} & \textbf{Analisis} \\
		\hline
		Cost solusi & Semua algoritma menghasilkan cost 1637. \\
		\hline
		Lava dan obstacle & Banyak jalur berisiko menyebabkan game over, sehingga validasi gerakan sangat penting. \\
		\hline
		Jumlah iterasi & UCS dan GBFS membutuhkan 11 iterasi, sedangkan A* dan WA* membutuhkan 12 iterasi. \\
		\hline
		Waktu eksekusi & UCS menjadi yang tercepat pada hasil percobaan ini. \\
		\hline
		Kesimpulan & Pada map ini, semua algoritma stabil dan menghasilkan kualitas solusi yang sama. \\
		\hline
	\end{tabular}
\end{table}


\section{Analisis Per Algoritma}

\subsection{Uniform Cost Search}

UCS selalu memilih state berdasarkan cost aktual terkecil. 
Berdasarkan hasil percobaan, UCS konsisten menghasilkan cost solusi terbaik atau sama dengan algoritma terbaik lainnya pada seluruh test case.

\begin{table}[H]
	\centering
	\caption{Analisis performa UCS}
	\begin{tabular}{|p{4cm}|p{10cm}|}
		\hline
		\rowcolor{green!15}
		\textbf{Aspek} & \textbf{Analisis} \\
		\hline
		Optimalitas & UCS paling aman untuk mencari total cost minimum karena selalu memperluas state dengan cost aktual terkecil. \\
		\hline
		Map kecil & Pada Test 8, UCS menghasilkan cost optimal dengan iterasi sedikit. \\
		\hline
		Map besar & Pada Test 17, UCS membutuhkan lebih banyak iterasi dibanding GBFS, tetapi tetap menghasilkan cost terbaik. \\
		\hline
		Map banyak tembok & Pada Test 20, UCS menghasilkan cost 810, jauh lebih baik daripada GBFS yang menghasilkan cost 6174. \\
		\hline
		Kelebihan & Sangat cocok untuk map dengan cost tile bervariasi. \\
		\hline
		Kekurangan & Tidak menggunakan informasi arah tujuan, sehingga pada beberapa kasus jumlah iterasi lebih besar. \\
		\hline
	\end{tabular}
\end{table}

\subsection{Greedy Best-First Search}

GBFS memilih state berdasarkan nilai heuristik terkecil. 
Hasil percobaan menunjukkan bahwa GBFS sering memiliki jumlah iterasi paling sedikit, tetapi tidak selalu menghasilkan cost terbaik.

\begin{table}[H]
	\centering
	\caption{Analisis performa GBFS}
	\begin{tabular}{|p{4cm}|p{10cm}|}
		\hline
		\rowcolor{orange!20}
		\textbf{Aspek} & \textbf{Analisis} \\
		\hline
		Kecepatan & GBFS sering menjadi algoritma dengan waktu eksekusi paling cepat. \\
		\hline
		Jumlah iterasi & Pada Test 8, Test 17, Test 20, dan Test 24, GBFS memiliki iterasi lebih sedikit dibanding algoritma lain. \\
		\hline
		Risiko optimalitas & Pada Test 20, GBFS menghasilkan cost 6174, jauh lebih besar daripada cost terbaik 810. \\
		\hline
		Pengaruh heuristik & Performa GBFS sangat bergantung pada kualitas heuristik yang digunakan. \\
		\hline
		Kelebihan & Sangat efisien untuk menemukan solusi dengan cepat. \\
		\hline
		Kekurangan & Tidak aman digunakan jika tujuan utama adalah cost minimum. \\
		\hline
	\end{tabular}
\end{table}

\subsection{A*}

A* menggabungkan cost aktual dan heuristik melalui fungsi $f(n)=g(n)+h(n)$. 
Berdasarkan hasil percobaan, A* menghasilkan cost terbaik pada seluruh test case, termasuk pada Test 20 yang menunjukkan kelemahan GBFS.

\begin{table}[H]
	\centering
	\caption{Analisis performa A*}
	\begin{tabular}{|p{4cm}|p{10cm}|}
		\hline
		\rowcolor{red!15}
		\textbf{Aspek} & \textbf{Analisis} \\
		\hline
		Keseimbangan & A* menyeimbangkan cost aktual dan estimasi menuju target. \\
		\hline
		Kualitas solusi & A* menghasilkan cost yang sama dengan UCS pada seluruh test case. \\
		\hline
		Jumlah iterasi & A* dapat memiliki iterasi lebih besar dari GBFS, tetapi lebih aman dari sisi cost. \\
		\hline
		Map banyak tembok & Pada Test 20, A* tetap menghasilkan cost 810, sedangkan GBFS menghasilkan cost 6174. \\
		\hline
		Stabilitas & A* stabil pada berbagai bentuk map karena tetap mempertimbangkan cost aktual. \\
		\hline
		Kesimpulan & A* merupakan algoritma paling seimbang antara optimalitas dan efisiensi pencarian. \\
		\hline
	\end{tabular}
\end{table}

\subsection{Weighted A*}

Weighted A* menggunakan fungsi evaluasi $f(n)=g(n)+2h(n)$, sehingga heuristik memiliki pengaruh lebih besar daripada A* biasa. 
Tujuannya adalah mempercepat pencarian dengan membuat algoritma lebih agresif menuju target.

\begin{table}[H]
	\centering
	\caption{Analisis performa Weighted A*}
	\begin{tabular}{|p{4cm}|p{10cm}|}
		\hline
		\rowcolor{purple!15}
		\textbf{Aspek} & \textbf{Analisis} \\
		\hline
		Pengaruh bobot & Bobot $w=2$ membuat algoritma lebih mengikuti heuristik dibanding A* biasa. \\
		\hline
		Kualitas solusi & Pada percobaan ini, Weighted A* menghasilkan cost yang sama dengan A* pada seluruh test case. \\
		\hline
		Jumlah iterasi & Jumlah iterasi Weighted A* umumnya sama dengan A* pada data uji yang digunakan. \\
		\hline
		Waktu eksekusi & Pada beberapa kasus, Weighted A* lebih cepat daripada A*, misalnya Test 20 dan Test 25. \\
		\hline
		Risiko & Secara teori, Weighted A* tidak selalu menjamin solusi optimal karena heuristik diperbesar. \\
		\hline
		Kesimpulan & Weighted A* cocok sebagai algoritma bonus untuk menguji trade-off antara kecepatan dan optimalitas. \\
		\hline
	\end{tabular}
\end{table}


\section{Analisis Pengaruh Heuristik}

Heuristik yang digunakan pada percobaan adalah H1, H2, dan H3. 
Berdasarkan implementasi program, H1 merepresentasikan Manhattan menuju goal, H2 mengarah ke angka berikutnya atau goal, sedangkan H3 memperkirakan beberapa target yang belum dikumpulkan.

\begin{longtable}{|p{2cm}|p{4cm}|p{7.5cm}|}
	\caption{Karakteristik heuristik} \\
	\hline
	\rowcolor{orange!20}
	\textbf{Heuristik} & \textbf{Nama} & \textbf{Karakteristik} \\
	\hline
	\endfirsthead
	
	\hline
	\rowcolor{orange!20}
	\textbf{Heuristik} & \textbf{Nama} & \textbf{Karakteristik} \\
	\hline
	\endhead
	
	H1 & Basic / Manhattan & Menghitung jarak Manhattan langsung menuju goal. Cepat dihitung, tetapi tidak memperhatikan angka berikutnya. \\
	\hline
	H2 & Custom & Menghitung jarak ke angka berikutnya yang harus dikumpulkan, lalu menuju goal jika semua angka selesai. \\
	\hline
	H3 & Advanced & Mengestimasi sisa perjalanan ke target-target yang belum dikumpulkan dan goal dengan pendekatan greedy. \\
	\hline
\end{longtable}

\begin{table}[H]
	\centering
	\caption{Analisis pengaruh heuristik terhadap algoritma}
	\begin{tabular}{|p{4cm}|p{10cm}|}
		\hline
		\rowcolor{orange!20}
		\textbf{Aspek} & \textbf{Analisis} \\
		\hline
		Pengaruh pada GBFS & Sangat besar, karena GBFS hanya menggunakan $h(n)$ sebagai prioritas. Pada Test 17 dan Test 24, heuristik Custom dan Advanced menurunkan jumlah iterasi secara signifikan. \\
		\hline
		Pengaruh pada A* & Cukup besar, tetapi tidak sedominan GBFS karena A* tetap mempertimbangkan $g(n)$. \\
		\hline
		Pengaruh pada Weighted A* & Lebih besar dibanding A* karena nilai heuristik dikalikan bobot $2$. \\
		\hline
		Heuristik paling cepat & Pada beberapa test case, Custom menghasilkan waktu terbaik, misalnya Test 17 dan Test 24 pada GBFS. \\
		\hline
		Heuristik paling stabil & Custom cukup stabil untuk map dengan angka karena targetnya mengikuti angka berikutnya. \\
		\hline
		Risiko heuristik & Heuristik yang tidak memperhitungkan obstacle, lava, dan sliding dapat mengarahkan algoritma ke jalur yang tampak dekat tetapi tidak optimal. \\
		\hline
	\end{tabular}
\end{table}


\section{Perbandingan Teoritis dan Eksperimental}

Secara teori, UCS menjamin cost minimum untuk cost positif, GBFS cepat tetapi tidak optimal, A* menyeimbangkan cost dan heuristik, sedangkan Weighted A* memperbesar pengaruh heuristik untuk mempercepat pencarian. 
Hasil percobaan secara umum sesuai dengan teori tersebut.

\begin{table}[H]
	\centering
	\caption{Perbandingan teori dan hasil eksperimen}
	\begin{tabular}{|p{3cm}|p{5cm}|p{6cm}|}
		\hline
		\rowcolor{gray!20}
		\textbf{Algoritma} & \textbf{Ekspektasi Teoritis} & \textbf{Hasil Eksperimen} \\
		\hline
		UCS & Optimal tetapi dapat mengeksplorasi banyak state. & Menghasilkan cost terbaik pada semua test case, tetapi tidak selalu memiliki iterasi paling kecil. \\
		\hline
		GBFS & Cepat, tetapi tidak selalu optimal. & Sering memiliki iterasi paling kecil, tetapi pada Test 20 menghasilkan cost jauh lebih besar. \\
		\hline
		A* & Seimbang antara cost aktual dan heuristik. & Menghasilkan cost terbaik dan lebih stabil daripada GBFS. \\
		\hline
		Weighted A* & Lebih agresif mengikuti heuristik dan dapat lebih cepat, tetapi tidak selalu optimal. & Pada percobaan ini cost tetap sama dengan A*, tetapi secara teori masih memiliki risiko menghasilkan cost lebih besar. \\
		\hline
	\end{tabular}
\end{table}

\begin{tcolorbox}[
	colback=gray!5,
	colframe=black!60,
	title=\textbf{Faktor Penyebab Perbedaan Hasil},
	arc=2mm,
	boxrule=0.8pt
]
Perbedaan hasil antar algoritma dipengaruhi oleh beberapa faktor, yaitu bentuk map, jumlah tembok, keberadaan angka, variasi cost tile, kualitas heuristik, banyaknya gerakan yang menyebabkan game over, dan ukuran papan. 
Pada map dengan banyak rintangan, heuristik jarak dapat menjadi kurang akurat karena jalur yang tampak dekat belum tentu dapat dicapai dengan cost rendah.
\end{tcolorbox}


\section{Kesimpulan Hasil Percobaan}

Berdasarkan seluruh hasil percobaan, dapat disimpulkan bahwa setiap algoritma memiliki keunggulan dan kelemahan masing-masing.

\begin{table}[H]
	\centering
	\caption{Kesimpulan hasil percobaan}
	\begin{tabular}{|p{5cm}|p{9cm}|}
		\hline
		\rowcolor{green!15}
		\textbf{Kategori} & \textbf{Kesimpulan} \\
		\hline
		Algoritma terbaik dari sisi optimalitas & UCS dan A* karena selalu menghasilkan cost terbaik pada seluruh test case. \\
		\hline
		Algoritma terbaik dari sisi waktu & GBFS sering menjadi yang tercepat, terutama pada Test 8, Test 17, Test 20, dan Test 24. \\
		\hline
		Algoritma terbaik dari sisi iterasi & GBFS karena sering meninjau state paling sedikit. \\
		\hline
		Algoritma terbaik untuk map kecil & GBFS cukup unggul karena dapat menemukan solusi dengan iterasi sedikit. \\
		\hline
		Algoritma terbaik untuk map banyak angka & GBFS dengan heuristik Custom atau Advanced efektif mengurangi iterasi, tetapi A* lebih aman untuk menjaga kualitas cost. \\
		\hline
		Algoritma terbaik untuk map banyak tembok & UCS dan A* lebih aman karena mempertimbangkan cost aktual. \\
		\hline
		Heuristik terbaik secara umum & Custom cukup baik karena mempertimbangkan angka berikutnya, sehingga lebih sesuai dengan constraint puzzle. \\
		\hline
		Rekomendasi algoritma utama & A* direkomendasikan sebagai algoritma utama karena memberikan keseimbangan antara kualitas cost solusi dan efisiensi pencarian. \\
		\hline
	\end{tabular}
\end{table}

\begin{tcolorbox}[
	colback=green!5,
	colframe=green!50!black,
	title=\textbf{Kesimpulan Akhir},
	arc=2mm,
	boxrule=0.8pt
]
UCS paling aman untuk memperoleh solusi dengan cost minimum, tetapi dapat membutuhkan eksplorasi lebih besar. 
GBFS paling cepat dari sisi iterasi dan waktu, tetapi berisiko menghasilkan solusi mahal seperti pada Test 20. 
A* menjadi pilihan paling seimbang karena tetap mempertimbangkan cost aktual dan heuristik. 
Weighted A* berguna sebagai algoritma bonus untuk mengevaluasi trade-off antara kecepatan dan optimalitas, terutama ketika heuristik ingin diberi pengaruh lebih besar.
\end{tcolorbox}